% Здесь пишется содержание приложений
%
\chapter{Листинг программного кода лабораторной работы №1 первой части}\label{app:A}\vspace{\baselineskip}

%Здесь можно привести полный листинг кода программы или модуля.

\begin{code}
import os
from pyspark.sql import SparkSession, DataFrame
from pyspark import SparkConf
from pyspark.sql.functions import (
    col,
    trim,
    unix_timestamp
)

def create_spark_configuration() -> SparkConf:
    """
    Создает и конфигурирует экземпляр SparkConf для приложения Spark.

    Returns:
        SparkConf: Настроенный экземпляр SparkConf.
    """
    # Получаем имя пользователя
    user_name = os.getenv("USER")
    
    conf = SparkConf()
    conf.setAppName("lab 1 US Accidents")
    conf.setMaster("yarn")
    conf.set("spark.submit.deployMode", "client")
    conf.set("spark.executor.memory", "12g")
    conf.set("spark.executor.cores", "8")
    conf.set("spark.executor.instances", "2")
    conf.set("spark.driver.memory", "4g")
    conf.set("spark.driver.cores", "2")
    conf.set("spark.jars.packages", "org.apache.iceberg:iceberg-spark-runtime-3.5_2.12:1.6.0")
    conf.set("spark.sql.extensions", "org.apache.iceberg.spark.extensions.IcebergSparkSessionExtensions")
    conf.set("spark.sql.catalog.spark_catalog", "org.apache.iceberg.spark.SparkCatalog")
    conf.set("spark.sql.catalog.spark_catalog.type", "hadoop")
    conf.set("spark.sql.catalog.spark_catalog.warehouse", f"hdfs:///user/{user_name}/warehouse")
    conf.set("spark.sql.catalog.spark_catalog.io-impl", "org.apache.iceberg.hadoop.HadoopFileIO")

    return conf

conf = create_spark_configuration()

spark = SparkSession.builder.config(conf=conf).getOrCreate()
spark

path = "hdfs:///user/user7/otpushchennikov/data/US_Accidents_March23.csv"

df = (spark.read.format("csv")
      .option("header", "true")
      .load(path)
)

df.show(3, truncate=25)

df = df.select(
    "ID", "Severity", "Start_Time", "End_Time", "Start_Lat", "Start_Lng", "Distance(mi)",
    "Street", "City", "County", "State", "Zipcode", "Timezone", "Airport_Code",
    "Weather_Timestamp", "Temperature(F)", "Wind_Chill(F)", "Humidity(%)", "Pressure(in)",
    "Visibility(mi)", "Wind_Direction", "Wind_Speed(mph)", "Precipitation(in)",
    "Weather_Condition", "Bump", "Crossing", "Give_Way", "Junction", "No_Exit", "Railway",
    "Roundabout", "Station", "Stop", "Traffic_Signal", "Turning_Loop", "Sunrise_Sunset",
    "Civil_Twilight", "Nautical_Twilight", "Astronomical_Twilight"
)

df.show(3, truncate=25)

df.printSchema()

BOOLEAN_COLUMNS = [
    "Bump", "Crossing", "Give_Way", "Junction", "No_Exit", "Railway", "Roundabout",
    "Station", "Stop", "Traffic_Signal", "Turning_Loop"
]

FLOAT_COLUMNS = [
    "Distance(mi)", "Temperature(F)", "Wind_Chill(F)", "Humidity(%)", "Pressure(in)",
    "Visibility(mi)", "Wind_Speed(mph)", "Precipitation(in)"
]

def transform_dataframe(data: DataFrame) -> DataFrame:
    """
    Преобразует столбцы DataFrame в указанные типы данных и
    выполняет необходимые преобразования.

    Args:
        data (DataFrame): Исходный DataFrame.

    Returns:
        DataFrame: Преобразованный DataFrame.
    """
    # Целевой признак и координаты
    data = data.withColumn("Severity",
                           col("Severity").cast("Integer"))
    data = data.withColumn("Start_Lat",
                           col("Start_Lat").cast("Double"))
    data = data.withColumn("Start_Lng",
                           col("Start_Lng").cast("Double"))

    # Временные метки
    data = data.withColumn("Start_Time",
                           col("Start_Time").cast("Timestamp"))
    data = data.withColumn("End_Time",
                           col("End_Time").cast("Timestamp"))
    data = data.withColumn("Weather_Timestamp",
                           col("Weather_Timestamp").cast("Timestamp"))

    # Производный признак: длительность происшествия в минутах
    data = data.withColumn("Duration_min",
                           ((unix_timestamp(col("End_Time")) - unix_timestamp(col("Start_Time"))) / 60).cast("Float"))

    # Числовые метеопризнаки: удалять единицы измерения не требуется,
    # они вынесены в заголовки столбцов
    for name in FLOAT_COLUMNS:
        data = data.withColumn(name, col(name).cast("Float"))

    # Булевы признаки инфраструктуры: в исходных данных это строки
    # 'True'/'False'
    for name in BOOLEAN_COLUMNS:
        data = data.withColumn(name, col(name).cast("Boolean"))

    # Идентификатор, почтовый индекс и адресные поля оставляем строками:
    # индекс -- не число, а код (ведущие нули значимы)
    data = data.withColumn("Zipcode",
                           trim(col("Zipcode")))

    return data

df = transform_dataframe(df)

df.show(3, truncate=25))

df.printSchema()

df.rdd.getNumPartitions()

database_name = "otpushchennikov_database"
table_name = "otpushchennikov_us_accidents"

create_database_sql = f"""
CREATE DATABASE IF NOT EXISTS spark_catalog.{database_name}
"""

spark.sql(create_database_sql)

spark.catalog.setCurrentDatabase(database_name)

# Сохранение DataFrame в виде таблицы
df.writeTo(table_name).using("iceberg").create()

for table in spark.catalog.listTables():
    print(table.name)

# spark.sql("DROP TABLE spark_catalog.otpushchennikov_database.otpushchennikov_us_accidents")
# spark.sql("DROP DATABASE spark_catalog.otpushchennikov_database")

spark.stop()
\end{code}

\chapter{Листинг программного кода лабораторной работы №1 второй части}\label{app:B}\vspace{\baselineskip}

\begin{code}
import os
from pyspark.sql import SparkSession, DataFrame
from pyspark import SparkConf
from pyspark.sql.functions import (
    col,
    lit,
    when,
    count,
    mean,
    sum,
    desc,
    floor,
    corr,
    expr,
    trim,
    lower,
    length,
    substring,
    month,
    dayofweek,
    hour
)
import pandas as pd
import seaborn as sns
import matplotlib.pyplot as plt

def create_spark_configuration() -> SparkConf:
    """
    Создает и конфигурирует экземпляр SparkConf для приложения Spark.

    Returns:
        SparkConf: Настроенный экземпляр SparkConf.
    """
    # Получаем имя пользователя
    user_name = os.getenv("USER")
    
    conf = SparkConf()
    conf.setAppName("lab 1 US Accidents. Part 2")
    conf.setMaster("yarn")
    conf.set("spark.submit.deployMode", "client")
    conf.set("spark.executor.memory", "12g")
    conf.set("spark.executor.cores", "8")
    conf.set("spark.executor.instances", "2")
    conf.set("spark.driver.memory", "4g")
    conf.set("spark.driver.cores", "2")
    conf.set("spark.jars.packages", "org.apache.iceberg:iceberg-spark-runtime-3.5_2.12:1.6.0")
    conf.set("spark.sql.extensions", "org.apache.iceberg.spark.extensions.IcebergSparkSessionExtensions")
    conf.set("spark.sql.catalog.spark_catalog", "org.apache.iceberg.spark.SparkCatalog")
    conf.set("spark.sql.catalog.spark_catalog.type", "hadoop")
    conf.set("spark.sql.catalog.spark_catalog.warehouse", f"hdfs:///user/{user_name}/warehouse")
    conf.set("spark.sql.catalog.spark_catalog.io-impl", "org.apache.iceberg.hadoop.HadoopFileIO")

    return conf

conf = create_spark_configuration()

spark = SparkSession.builder.config(conf=conf).getOrCreate()
spark

database_name = "otpushchennikov_database"
table_name = "otpushchennikov_us_accidents"

spark.catalog.setCurrentDatabase(database_name)

df = spark.table(table_name)

df.show(3, truncate=25)

df.printSchema()

df.count()

df = df.cache()

def missing_overview(data: DataFrame) -> pd.DataFrame:
    """
    Подсчет количества и доли пропущенных значений для всех столбцов
    датафрейма за один проход по данным.

    Args:
        data (DataFrame): DataFrame, содержащий данные.

    Returns:
        pd.DataFrame: Таблица со столбцами `column`, `nulls`, `share`,
            отсортированная по убыванию доли пропусков.
    """
    total = data.count()
    counters = [count(when(col(name).isNull(), lit(True))).alias(str(index))
                for index, name in enumerate(data.columns)]
    counts = data.agg(*counters).collect()[0].asDict()

    rows = [
        (name, int(counts[str(index)]), 100.0 * int(counts[str(index)]) / total)
        for index, name in enumerate(data.columns)
    ]
    overview = pd.DataFrame(rows, columns=["column", "nulls", "share"])
    return overview.sort_values("share", ascending=False).reset_index(drop=True)

missing = missing_overview(df)
missing[missing["share"] > 0].to_string(index=False)

plot_data = missing[missing["share"] > 0]
plt.figure(figsize=(10, 6))
sns.barplot(x="share", y="column", data=plot_data)
plt.title("Доля пропущенных значений по столбцам")
plt.xlabel("share, %")
plt.ylabel("")
plt.grid(True)
plt.show()

(df.select("ID")
  .filter(~col("ID").rlike("^[A-Z]+-[0-9]+$"))
  .count())

(df.groupBy("ID")
   .count()
   .where("count > 1")
   .count())

df.select("ID").distinct().count()

def count_nulls(data: DataFrame,
                column_name: str) -> None:
    """
    Подсчет количества null и not null значений в указанном столбце.

    Args:
        data (DataFrame): DataFrame, содержащий данные.
        column_name (str): Имя столбца для подсчета null и not null значений.
    
    Returns:
        None
    """
    # Подсчет количества null значений в указанном столбце
    null_counts = data.select(
        sum(col(column_name).isNull().cast("int"))
    ).collect()[0][0]

    # Подсчет количества not null значений в указанном столбце
    not_null_counts = data.select(
        sum(col(column_name).isNotNull().cast("int"))
    ).collect()[0][0]

    # Вывод результатов
    print(f"Число колонок с NULL: {null_counts} "
          f"({100 * null_counts / (null_counts + not_null_counts):.2f}%)")

count_nulls(data=df, column_name="Severity")

def plot_cat_distribution(data: DataFrame,
                          column_name: str,
                          top_n: int = 20) -> None:
    """
    Построение распределения категориального признака.

    Args:
        data (DataFrame): DataFrame, содержащий данные.
        column_name (str): Имя столбца для группировки.
        top_n (int): Количество топ-значений для отображения.
    
    Returns:
        None
    """
    # Группировка данных по столбцу и подсчет количества
    categories = (
        data
        .groupBy(column_name)
        .count()
        .orderBy("count", ascending=False)
    )
    
    print(f"Количество категорий признака {column_name}: {categories.count()}")

    categories = (
        categories
        .limit(top_n)
        .toPandas()
    )
    
    # Визуализация с использованием Seaborn
    plt.figure(figsize=(10, 6))
    sns.barplot(x=column_name, y="count", data=categories)
    plt.title(f"Barplot of \"{column_name}\" counts")
    plt.xlabel(column_name)
    plt.ylabel("Count")
    plt.xticks(rotation=45)
    plt.show()

plot_cat_distribution(data=df, column_name="Severity")

count_nulls(data=df, column_name="Duration_min")

def plot_boxplots(data: DataFrame,
                  columns: list[str],
                  sample_fraction: float = 0.1) -> None:
    """
    Построение boxplot для нескольких столбцов в PySpark DataFrame.

    Args:
        data (DataFrame): DataFrame, содержащий данные.
        columns (list of str): Список имен столбцов для построения boxplot.
        sample_fraction (float): Доля данных для семплирования выбросов.
    
    Returns:
        None
    """
    box_data = []

    for column in columns:
        # Вычисление квантилей
        quantiles = data.approxQuantile(column, [0.25, 0.5, 0.75], 0.01)
        q1, median, q3 = quantiles

        # Вычисление IQR и границ усов
        iqr = q3 - q1
        lower_bound = q1 - 1.5 * iqr
        upper_bound = q3 + 1.5 * iqr

        # Вычисление минимального и максимального значений
        min_value = data.agg({column: "min"}).collect()[0][0]
        mean_value = data.agg({column: "mean"}).collect()[0][0]
        std_value = data.agg({column: "stddev"}).collect()[0][0]
        max_value = data.agg({column: "max"}).collect()[0][0]

        # Фильтрация выбросов
        outliers_df = data.filter((col(column) < lower_bound) | (col(column) > upper_bound))
        outliers_count = outliers_df.count()

        # Ограничение усов минимальным и максимальным значениями
        lower_bound = max(lower_bound, min_value)
        upper_bound = min(upper_bound, max_value)

        # Семплирование выбросов
        outliers = []
        if outliers_count > 0:
            sampled_outliers_df = outliers_df.sample(sample_fraction)
            outliers = (
                sampled_outliers_df
                .select(column)
                .limit(1000)
                .collect()
            )
            outliers = [row[column] for row in outliers]

        # Вывод статистических характеристик
        print(f"{column}: {outliers_count} значений ("
              f"{100 * outliers_count / data.count():.2f}%) вне границ "
              f"[{lower_bound:.2f}, {upper_bound:.2f}]")
        print(f"  Минимальное значение:          {min_value:.2f}")
        print(f"  Среднее значение:              {mean_value:.2f}")
        print(f"  Среднеквадратичное отклонение: {std_value:.2f}")
        print(f"  Первый квартиль:               {q1:.2f}")
        print(f"  Медиана:                       {median:.2f}")
        print(f"  Третий квартиль:               {q3:.2f}")
        print(f"  Максимальное значение:         {max_value:.2f}")

        # Подготовка данных для axes.bxp
        box_data.append({
            'whislo': lower_bound,  # Нижняя граница усов
            'q1': q1,               # Первый квартиль
            'med': median,          # Медиана
            'q3': q3,               # Третий квартиль
            'whishi': upper_bound,  # Верхняя граница усов
            'fliers': outliers      # Выбросы
        })

    # Построение boxplot
    fig, ax = plt.subplots(figsize=(20, 6))
    ax.bxp(box_data,
           vert=False,
           positions=range(1, len(columns) + 1), widths=0.5)
    ax.set_yticks(range(1, len(columns) + 1))
    ax.set_yticklabels(columns)
    ax.set_xlabel('Value')
    ax.set_title('Boxplots')
    ax.grid(True)
    plt.show()

plot_boxplots(data=df, columns=["Duration_min"])

def plot_quant_distribution(data: DataFrame,
                            column: str,
                            num_bins: int = 200) -> None:
    """
    Построение гистограммы для количественной переменной с
    использованием PySpark и Seaborn.

    Args:
        data (DataFrame): DataFrame с данными.
        column (str): Название колонки с количественной переменной.
        num_bins (int): Количество бинов для гистограммы.

    Returns:
        None
    """
    data = data.filter(col(column).isNotNull())

    # Находим минимальное и максимальное значения колонки
    min_value = data.agg({column: "min"}).collect()[0][0]
    max_value = data.agg({column: "max"}).collect()[0][0]

    # Размер бина
    bin_size = (max_value - min_value) / num_bins

    # Добавляем колонку с номером бина
    data = data.withColumn(
        "bin",
        floor((col(column) - min_value) / bin_size)
    )

    # Группируем по номеру бина и считаем количество строк в каждом бине
    bin_counts = data.groupBy("bin").count()

    # Преобразуем результат в Pandas DataFrame для построения гистограммы
    bin_counts_pd = bin_counts.limit(1000).toPandas()
    
    # Создаем массив границ бинов
    bin_edges = [min_value + i * bin_size for i in range(num_bins + 2)]
    
    # Преобразуем номера бинов в центры бинов
    bin_centers = [
        (bin_edges[i] + bin_edges[i + 1]) / 2 for i in range(num_bins + 1)
    ]
    
    # Добавляем центры бинов в Pandas DataFrame
    bin_counts_pd['bin_center'] = bin_counts_pd['bin'].apply(
        lambda x: bin_centers[int(x)]
    )
    
    # Построение гистограммы с использованием Seaborn
    plt.figure(figsize=(20, 6))
    sns.histplot(data=bin_counts_pd, x="bin_center",
                 weights="count", kde=True, bins=num_bins + 1)
    plt.xlabel("Value")
    plt.ylabel("Count")
    plt.title(f"Распределение количественного признака \"{column}\"")
    plt.grid(True)
    plt.show()

plot_quant_distribution(data=df, column="Duration_min")

total = df.count()
for limit_minutes in (1440, 2880, 10080):
    n = df.filter(col("Duration_min") > limit_minutes).count()
    print(f"Длительность больше {limit_minutes} мин: {n} записей "
          f"({100.0 * n / total:.3f}%)")

DURATION_CAP = 2880.0
df = df.withColumn(
    "Duration_min",
    when(col("Duration_min") > DURATION_CAP, DURATION_CAP).otherwise(col("Duration_min"))
)
plot_quant_distribution(data=df, column="Duration_min")

count_nulls(data=df, column_name="Distance(mi)")

plot_boxplots(data=df, columns=["Distance(mi)"])

(df.filter(col("Distance(mi)") > 60).count(),
df.filter(col("Distance(mi)") > 100).count())

DISTANCE_CAP = 60.0
df = df.withColumn(
    "Distance(mi)",
    when(col("Distance(mi)") > DISTANCE_CAP, DISTANCE_CAP).otherwise(col("Distance(mi)")
    )
)
plot_quant_distribution(data=df, column="Distance(mi)")

WEATHER_COLUMNS = [
    "Temperature(F)", "Wind_Chill(F)", "Humidity(%)", "Pressure(in)",
    "Visibility(mi)", "Wind_Speed(mph)", "Precipitation(in)"
]
df.select(WEATHER_COLUMNS).describe().show(truncate=8)

PHYSICAL_BOUNDS = {
    "Temperature(F)": (-60.0, 130.0),   # правдоподобные температуры воздуха
    "Wind_Chill(F)": (-60.0, 130.0),
    "Humidity(%)": (0.0, 100.0),        # физический предел шкалы влажности
    "Pressure(in)": (25.0, 32.0),       # шкала inHg для наземных наблюдений
    "Visibility(mi)": (0.0, 90.0),
    "Wind_Speed(mph)": (0.0, 120.0),
    "Precipitation(in)": (0.0, 3.0),    # осадки за одно наблюдение
}

for name, (low, high) in PHYSICAL_BOUNDS.items():
    n = df.filter((col(name) < low) | (col(name) > high)).count()
    print(f"{name}: значений вне [{low}, {high}] -- {n} "
          f"({100.0 * n / total:.3f}%)")

def clip_out_of_bounds(data: DataFrame, bounds: dict) -> DataFrame:
    """
    Заменяет значения столбцов, выходящие за физически допустимые
    границы, на пропуски.

    Args:
        data (DataFrame): DataFrame, содержащий данные.
        bounds (dict): Сопоставление имени столбца и кортежа (low, high).

    Returns:
        DataFrame: Датафрейм с очищенными значениями.
    """
    for name, (low, high) in bounds.items():
        data = data.withColumn(
            name,
            when((col(name) >= low) & (col(name) <= high), col(name))
            .otherwise(lit(None).cast(data.schema[name].dataType))
        )
    return data


df = clip_out_of_bounds(df, PHYSICAL_BOUNDS)

MEDIAN_FILL = ["Temperature(F)", "Wind_Chill(F)", "Humidity(%)", "Pressure(in)",
               "Visibility(mi)", "Wind_Speed(mph)"]

for name in MEDIAN_FILL:
    median_value = df.approxQuantile(name, [0.5], 0.0001)[0]
    print(f"Медиана {name}: {median_value:.2f}")
    df = df.fillna({name: median_value})

df = df.fillna({"Precipitation(in)": 0.0})

after = missing_overview(df)
after[after["column"].isin(WEATHER_COLUMNS)].to_string(index=False)

plot_quant_distribution(data=df, column="Temperature(F)")

plot_quant_distribution(data=df, column="Wind_Speed(mph)")

CARDINALITY_COLUMNS = ["State", "Timezone", "Weather_Condition", "County",
                          "Airport_Code", "City", "Street", "Zipcode",
                          "Wind_Direction"]

for name in CARDINALITY_COLUMNS:
    n = df.filter(col(name).isNotNull()).select(name).distinct().count()
    print(f"{name}: {n} различных значений")

(df.filter(col("Zipcode").isNull()).count(),
 df.filter(col("Zipcode").contains("-")).count(),
 df.filter(length(col("Zipcode")) == 5).count())

df = df.withColumn("Zipcode", substring(col("Zipcode"), 1, 5))
df.filter(col("Zipcode").isNotNull()).select("Zipcode").distinct().count()

df = df.drop("Street")
df = df.fillna({"Airport_Code": "Unknown"})
df.show(3, truncate=25)

plot_cat_distribution(data=df, column_name="State", top_n=10)

df = df.fillna({"Weather_Condition": "Unknown"})
count_nulls(data=df, column_name="Weather_Condition")
plot_cat_distribution(data=df, column_name="Weather_Condition", top_n=15)

KEEP_CONDITIONS = [
    row[0] for row in
    df.groupBy("Weather_Condition").count()
      .filter(col("count") > 0.001 * total)
      .collect()
]
print(f"Категорий с долей более 0,1%: {len(KEEP_CONDITIONS)} из "
      f"{df.select('Weather_Condition').distinct().count()}")

df = df.withColumn(
    "Weather_Condition",
    when(col("Weather_Condition").isin(KEEP_CONDITIONS), col("Weather_Condition"))
    .otherwise(lit("Other"))
)
plot_cat_distribution(data=df, column_name="Weather_Condition", top_n=15)

print(f"Различных написаний: {df.select('Wind_Direction').distinct().count()}")
print(f"После trim(lower): {df.select(trim(lower(col('Wind_Direction')))).distinct().count()}")

(df.select(trim(lower(col("Wind_Direction"))).alias("wd"))
   .groupBy("wd")
   .count()
   .orderBy("count", ascending=False)
   .show(30, truncate=False))

WIND_SYNONYMS = {"north": "n", "south": "s", "east": "e", "west": "w",
                  "variable": "var"}

df = df.withColumn("Wind_Direction", trim(lower(col("Wind_Direction"))))
for long_name, short_name in WIND_SYNONYMS.items():
    df = df.withColumn(
        "Wind_Direction",
        when(col("Wind_Direction") == long_name, lit(short_name))
        .otherwise(col("Wind_Direction"))
    )

df = df.fillna({"Wind_Direction": "unknown"})
df.select("Wind_Direction").distinct().count()

plot_cat_distribution(data=df, column_name="Wind_Direction", top_n=12)

for name in ["Timezone", "Sunrise_Sunset", "Civil_Twilight", "Nautical_Twilight",
              "Astronomical_Twilight", "City", "County", "Zipcode"]:
    df = df.fillna({name: "Unknown"})

remaining = missing_overview(df)
remaining[remaining["share"] > 0].to_string(index=False)

df = df.drop("Weather_Timestamp")

BOOLEAN_COLUMNS = [
    "Bump", "Crossing", "Give_Way", "Junction", "No_Exit", "Railway", "Roundabout",
    "Station", "Stop", "Traffic_Signal", "Turning_Loop"
]

counters = [count(when(col(name), lit(True))).alias(str(index))
            for index, name in enumerate(BOOLEAN_COLUMNS)]
true_counts = df.agg(*counters).collect()[0].asDict()

shares = pd.DataFrame(
    [(name, int(true_counts[str(index)]), 100.0 * int(true_counts[str(index)]) / total)
     for index, name in enumerate(BOOLEAN_COLUMNS)],
    columns=["column", "true", "share"]
).sort_values("share", ascending=False)
shares.to_string(index=False)

LOW_VARIANCE = [name for name, share in zip(shares["column"], shares["share"]) if share < 0.1]
print("Удаляемые признаки:", LOW_VARIANCE)
df = df.drop(*LOW_VARIANCE)
df.show(3, truncate=25)

plot_cat_distribution(data=df, column_name="Traffic_Signal")

df = (df
      .withColumn("Month", month(col("Start_Time")))
      .withColumn("Day_Of_Week", dayofweek(col("Start_Time")))
      .withColumn("Hour", hour(col("Start_Time"))))
df.select("Start_Time", "Month", "Day_Of_Week", "Hour").show(5, truncate=False)

month_counts = df.groupBy("Month").count().orderBy("Month").toPandas()
plt.figure(figsize=(12, 5))
sns.barplot(x="Month", y="count", data=month_counts)
plt.title("Число происшествий по месяцам")
plt.xlabel("месяц")
plt.ylabel("count")
plt.grid(True)
plt.show()

day_counts = df.groupBy("Day_Of_Week").count().orderBy("Day_Of_Week").toPandas()
# В Spark нумерация дней недели начинается с воскресенья: 1 = Sun, ..., 7 = Sat
day_counts["day_name"] = ["Sun", "Mon", "Tue", "Wed", "Thu", "Fri", "Sat"]
plt.figure(figsize=(10, 5))
sns.barplot(x="day_name", y="count", data=day_counts)
plt.title("Число происшествий по дням недели")
plt.xlabel("")
plt.ylabel("count")
plt.grid(True)
plt.show()

hour_counts = df.groupBy("Hour").count().orderBy("Hour").toPandas()
plt.figure(figsize=(12, 5))
sns.lineplot(x="Hour", y="count", marker="o", data=hour_counts)
plt.title("Число происшествий по часам суток")
plt.xlabel("час")
plt.ylabel("count")
plt.grid(True)
plt.show()

def plot_grouped_boxplots(data: DataFrame,
                          y_column: str,
                          x_column: str,
                          top_n: int = 20) -> None:
    """
    Диаграмма «ящик с усами» для количественного столбца в разрезе
    категориального столбца.

    Args:
        data (DataFrame): DataFrame с данными.
        y_column (str): количественный столбец (ось значений).
        x_column (str): категориальный столбец (ось групп).
        top_n (int): число самых частых категорий для вывода.

    Returns:
        None
    """
    groups = (data.groupBy(x_column).count()
              .orderBy(desc("count"))
              .limit(top_n)
              .toPandas()[x_column].tolist())

    stats = (data.filter(col(x_column).isin(groups))
             .groupBy(x_column)
             .agg(
                 expr(f"percentile_approx(`{y_column}`, "
                      f"array(0.25, 0.5, 0.75), 1000) AS q"),
                 expr(f"min(`{y_column}`) AS min_value"),
                 expr(f"max(`{y_column}`) AS max_value"),
                 expr(f"count(`{y_column}`) AS n"),
             )
             .toPandas())

    stats["q1"] = stats["q"].apply(lambda value: value[0])
    stats["median"] = stats["q"].apply(lambda value: value[1])
    stats["q3"] = stats["q"].apply(lambda value: value[2])
    stats["iqr"] = stats["q3"] - stats["q1"]
    stats["lower"] = [max(q1 - 1.5 * spread, smallest)
                      for q1, spread, smallest in
                      zip(stats["q1"], stats["iqr"], stats["min_value"])]
    stats["upper"] = [min(q3 + 1.5 * spread, largest)
                     for q3, spread, largest in
                     zip(stats["q3"], stats["iqr"], stats["max_value"])]
    stats = stats.sort_values("median")

    print(f"{y_column} по {x_column}: "
          f"медианы групп от {stats['median'].min():.2f} "
          f"до {stats['median'].max():.2f}, "
          f"размах медиан {stats['median'].max() - stats['median'].min():.2f}, "
          f"средний межквартильный размах {stats['iqr'].mean():.2f}")
    print(stats[[x_column, "n", "q1", "median", "q3", "iqr"]]
          .to_string(index=False))

    boxes = [
        {
            "whislo": row.lower,
            "q1": row.q1,
            "med": row.median,
            "q3": row.q3,
            "whishi": row.upper,
            "fliers": [],
        }
        for row in stats.itertuples()
    ]

    fig, ax = plt.subplots(figsize=(20, 6))
    ax.bxp(boxes, positions=range(1, len(boxes) + 1), widths=0.6)
    ax.set_xticks(range(1, len(boxes) + 1))
    ax.set_xticklabels(stats[x_column].astype(str).tolist(), rotation=90)
    ax.set_xlabel(x_column)
    ax.set_ylabel(y_column)
    ax.set_title(f"{y_column} в разрезе {x_column} "
                 f"(точки выбросов не показываются)")
    ax.grid(True)
    plt.show()

plot_grouped_boxplots(data=df, y_column="Temperature(F)",
                      x_column="State", top_n=20)

plot_grouped_boxplots(data=df, y_column="Temperature(F)",
                      x_column="Month", top_n=12)

plot_grouped_boxplots(data=df, y_column="Wind_Speed(mph)",
                      x_column="Month", top_n=12)

plot_grouped_boxplots(data=df, y_column="Duration_min",
                      x_column="Severity", top_n=4)

df = df.cache()
df.count()

def compute_and_visualize_correlation_matrix(data: DataFrame,
                                             columns: list[str]) -> None:
    """
    Вычисляет и визуализирует корреляционную матрицу для указанных
    колонок в DataFrame PySpark.

    Args:
        data (DataFrame): DataFrame PySpark.
        columns (list[str]): Список колонок для вычисления корреляции.

    Returns:
        None
    """
    # Вычисление корреляционной матрицы
    corr_matrix = {c1: {c2: 1.0 for c2 in columns} for c1 in columns}
    for index, col1 in enumerate(columns):
        for col2 in columns[index + 1:]:
            corr_value = data.select(corr(col1, col2)).collect()[0][0]
            corr_matrix[col1][col2] = corr_value
            corr_matrix[col2][col1] = corr_value

    # Преобразование корреляционной матрицы в DataFrame Pandas для визуализации
    corr_matrix_pd = pd.DataFrame(corr_matrix)

    # Числа самых сильных связей проще прочитать в тексте, чем на тепловой карте
    pairs = pd.DataFrame(
        [(col1, col2, corr_matrix[col1][col2])
         for index, col1 in enumerate(columns) for col2 in columns[index + 1:]],
        columns=["column_1", "column_2", "corr"]
    )
    pairs["abs_corr"] = pairs["corr"].abs()
    print(pairs.sort_values("abs_corr", ascending=False)
          .head(8)
          .to_string(index=False))

    # Построение и визуализация корреляционной матрицы
    plt.figure(figsize=(10, 8))
    sns.heatmap(corr_matrix_pd, annot=True, fmt=".2f", cmap="coolwarm", linewidths=0.5)
    plt.title("Correlation Matrix")
    plt.show()

compute_and_visualize_correlation_matrix(
    data=df, columns=[
        "Severity", "Distance(mi)", "Duration_min", "Temperature(F)",
        "Wind_Chill(F)", "Humidity(%)", "Pressure(in)", "Visibility(mi)",
        "Wind_Speed(mph)", "Precipitation(in)"
    ]
)

df.count()

df.printSchema()

# Сохранение DataFrame в виде таблицы
df.writeTo("otpushchennikov_us_accidents_processed").using("iceberg").create()

for table in spark.catalog.listTables():
    print(table.name)

# spark.sql("DROP TABLE spark_catalog.otpushchennikov_database.otpushchennikov_us_accidents_processed")
# spark.sql("DROP TABLE spark_catalog.otpushchennikov_database.otpushchennikov_us_accidents")
# spark.sql("DROP DATABASE spark_catalog.otpushchennikov_database")

spark.catalog.clearCache()
spark.stop()

\end{code}

\chapter{Листинг программного кода лабораторной работы №2 первой части}\label{app:C}\vspace{\baselineskip}

\begin{code}
import os
import csv
import time
import numpy as np
import pandas as pd
from pyspark.sql import SparkSession, DataFrame
from pyspark import SparkConf
from pyspark.sql.functions import col, lit, coalesce, mean, sqrt, approx_count_distinct
from pyspark.sql.types import DoubleType
from pyspark.ml import Pipeline
from pyspark.ml.feature import (
    VectorAssembler,
    StringIndexer,
    OneHotEncoder,
    MinMaxScaler,
    Binarizer
)
from pyspark.ml.regression import LinearRegression
from pyspark.ml.evaluation import RegressionEvaluator
from pyspark.ml.tuning import CrossValidator, CrossValidatorModel, ParamGridBuilder
import seaborn as sns
import matplotlib.pyplot as plt

def create_spark_configuration() -> SparkConf:
    """
    Создает и конфигурирует экземпляр SparkConf для приложения Spark.

    Returns:
        SparkConf: Настроенный экземпляр SparkConf.
    """
    # Получаем имя пользователя
    user_name = os.getenv("USER")

    conf = SparkConf()
    conf.setAppName("lab 2 US Accidents. Linear Regression")
    conf.setMaster("yarn")
    conf.set("spark.submit.deployMode", "client")
    conf.set("spark.executor.memory", "12g")
    conf.set("spark.executor.cores", "8")
    conf.set("spark.executor.instances", "2")
    conf.set("spark.driver.memory", "4g")
    conf.set("spark.driver.cores", "2")
    conf.set("spark.jars.packages",
             "org.apache.iceberg:iceberg-spark-runtime-3.5_2.12:1.6.0")
    conf.set("spark.sql.extensions",
             "org.apache.iceberg.spark.extensions.IcebergSparkSessionExtensions")
    conf.set("spark.sql.catalog.spark_catalog",
             "org.apache.iceberg.spark.SparkCatalog")
    conf.set("spark.sql.catalog.spark_catalog.type", "hadoop")
    conf.set("spark.sql.catalog.spark_catalog.warehouse",
             f"hdfs:///user/{user_name}/warehouse")
    conf.set("spark.sql.catalog.spark_catalog.io-impl",
             "org.apache.iceberg.hadoop.HadoopFileIO")

    return conf

conf = create_spark_configuration()

spark = SparkSession.builder.config(conf=conf).getOrCreate()
spark

database_name = "otpushchennikov_database"

spark.catalog.setCurrentDatabase(database_name)

df = spark.table("otpushchennikov_us_accidents_processed")

df.show()

df.printSchema()

df.count()

binary_flags = [
    "Crossing", "Give_Way", "Junction", "No_Exit",
    "Railway", "Station", "Stop", "Traffic_Signal"
]
float_features = [
    "Distance(mi)", "Temperature(F)", "Wind_Chill(F)", "Humidity(%)",
    "Pressure(in)", "Visibility(mi)", "Wind_Speed(mph)", "Precipitation(in)",
    "Month", "Day_Of_Week", "Hour"
]

df = df.withColumn("Duration_min", col("Duration_min").cast(DoubleType()))
for name in float_features:
    df = df.withColumn(name, col(name).cast(DoubleType()))
for flag in binary_flags:
    df = df.withColumn(
        flag,
        coalesce(col(flag).cast(DoubleType()), lit(0.0))
    )

null_label = df.filter(col("Duration_min").isNull()).count()
print(f"Строк без значения Duration_min: {null_label}")

if null_label > 0:
    df = df.filter(col("Duration_min").isNotNull())

print(f"Строк в работе: {df.count()}")

df.select(col("Duration_min").alias("value")).summary("50%", "90%", "99%").show()

candidate_categorical = [
    "State", "Weather_Condition", "Wind_Direction", "Sunrise_Sunset",
    "Timezone", "Airport_Code", "County", "City", "Zipcode"
]

cardinality = df.agg(
    *[approx_count_distinct(col(name)).alias(name) for name in candidate_categorical]
).toPandas().T
cardinality.columns = ["число уровней"]
print(cardinality.sort_values("число уровней", ascending=False).to_string())

def save_sample_to_csv(data: DataFrame, file_path: str,
                       sample_size: int = 1000) -> DataFrame:
    """
    Сохраняет первые `sample_size` строк из DataFrame в CSV-файл на драйвере
    и возвращает DataFrame с оставшимися данными.

    Args:
        data (DataFrame): DataFrame, из которого нужно извлечь строки.
        file_path (str): Путь для сохранения CSV-файла.
        sample_size (int): Количество строк для сохранения
            (по умолчанию 1000).

    Returns:
        DataFrame: DataFrame с оставшимися данными.
    """
    # Определяем пропорции для разделения
    sample_fraction = sample_size / data.count()
    remaining_fraction = 1 - sample_fraction

    # Разделяем DataFrame на два непересекающихся набора данных
    sample_data, remaining_data = data.randomSplit(
        [sample_fraction, remaining_fraction]
    )

    # Сохраняем извлеченные строки в CSV-файл на драйвере
    try:
        with open(file_path, mode="w", newline="") as file:
            writer = csv.writer(file)

            # Записываем заголовок
            writer.writerow(data.columns)

            # Записываем строки
            for row in sample_data.take(sample_size):
                writer.writerow(row)
        print(f"Файл \"{file_path}\" с данными успешно создан.")

    except Exception as e:
        print(f"Ошибка при записи файла: {e}")

    return remaining_data

path = "us_accidents_sample.csv"

df = save_sample_to_csv(data=df, file_path=path, sample_size=1000)

df.count()

train_df, test_df = df.randomSplit([0.8, 0.2])

print(f"Train dataset size: {train_df.count()}")
print(f"Test  dataset size: {test_df.count()}")

categorical_features = ["State", "Weather_Condition",
                      "Wind_Direction", "Sunrise_Sunset"]
binary_features = [
    "Crossing", "Give_Way", "Junction", "No_Exit",
    "Railway", "Station", "Stop", "Traffic_Signal"
]
numeric_features = [
    "Distance(mi)", "Temperature(F)", "Wind_Chill(F)", "Humidity(%)",
    "Pressure(in)", "Visibility(mi)", "Wind_Speed(mph)",
    "Month", "Day_Of_Week", "Hour"
]
binarizable_feature = "Precipitation(in)"
label_col = "Duration_min"
threshold = 0.1

total_rows = df.count()
wet_rows = df.filter(col(binarizable_feature) > threshold).count()
print(f"Записей с осадками > {threshold} in: {wet_rows} "
      f"({100.0 * wet_rows / total_rows:.2f} %)")

def create_pipeline(categorical_features: list[str], numeric_features: list[str],
                    binary_features: list[str], binarized_col: str,
                    threshold: float, label_col: str, max_iter: int) -> Pipeline:
    """
    Создает конвейер подготовки признаков и модель линейной регрессии.

    Args:
        categorical_features: Список категориальных признаков.
        numeric_features: Список количественных признаков.
        binary_features: Список бинарных признаков.
        binarized_col: Название бинаризуемого признака.
        threshold: Порог бинаризации.
        label_col: Название столбца с целевой переменной.
        max_iter: Максимальное количество итераций для модели.

    Returns:
        Pipeline: Конвейер подготовки признаков и обучения модели.
    """
    # Формируем названия колонок для преобразованных признаков
    indexed_categorical_features = [f"{feature}_index" for feature in categorical_features]
    onehot_categorical_features = [f"{feature}_ohe" for feature in categorical_features]

    # Выполняем преобразования данных
    string_indexer = StringIndexer(inputCols=categorical_features,
                                   outputCols=indexed_categorical_features,
                                   handleInvalid="keep")
    onehot_encoder = OneHotEncoder(inputCols=indexed_categorical_features,
                                   outputCols=onehot_categorical_features,
                                   dropLast=True,
                                   handleInvalid="keep")
    vector_num_assembler = VectorAssembler(inputCols=numeric_features,
                                           outputCol="numeric_vector",
                                           handleInvalid="keep")
    numeric_scaler = MinMaxScaler(inputCol="numeric_vector",
                                  outputCol="numeric_vector_scaled")
    binarizer = Binarizer(inputCol=binarized_col,
                          outputCol="precipitation_binary",
                          threshold=threshold)
    vector_all_assembler = VectorAssembler(
        inputCols=(onehot_categorical_features + binary_features +
                   ["numeric_vector_scaled", "precipitation_binary"]),
        outputCol="features",
        handleInvalid="keep"
    )

    # Создаем модель линейной регрессии
    linear_regression = LinearRegression(featuresCol="features",
                                         labelCol=label_col,
                                         predictionCol="prediction",
                                         standardization=False,
                                         maxIter=max_iter)

    # Создаем конвейер
    pipeline = Pipeline(stages=[
        string_indexer, onehot_encoder, vector_num_assembler,
        numeric_scaler, binarizer, vector_all_assembler,
        linear_regression
    ])

    return pipeline

pipeline = create_pipeline(categorical_features=categorical_features,
                           numeric_features=numeric_features,
                           binary_features=binary_features,
                           binarized_col=binarizable_feature,
                           threshold=threshold,
                           label_col=label_col,
                           max_iter=15)

param_grid = ParamGridBuilder() \
    .addGrid(pipeline.getStages()[-1].regParam, [0.01, 0.1, 1.0]) \
    .addGrid(pipeline.getStages()[-1].elasticNetParam, [0.0, 0.5, 1.0]) \
    .build()
print(f"Число комбинаций параметров: {len(param_grid)}")

cv_evaluator = RegressionEvaluator(labelCol=label_col,
                                   predictionCol="prediction",
                                   metricName="rmse")

train_df = train_df.cache()
test_df = test_df.cache()
print(f"Train rows cached: {train_df.count()}")

t0 = time.time()
single_pipeline_model = pipeline.fit(train_df)
one_fit_seconds = time.time() - t0

n_fits = len(param_grid) * 5
print(f"Время одного обучения конвейера: {one_fit_seconds:.1f} сек")
print(f"Число обучений в кросс-валидации: {n_fits}")
print(f"Оценка времени подбора: {one_fit_seconds * n_fits / 60:.1f} мин")

cross_validator = CrossValidator(estimator=pipeline,
                                 estimatorParamMaps=param_grid,
                                 evaluator=cv_evaluator,
                                 numFolds=5)

cv_model = cross_validator.fit(train_df)

def get_best_model_params(cv_model: CrossValidatorModel) -> dict[str, float]:
    """
    Получает параметры лучшей модели из объекта CrossValidatorModel.

    Args:
        cv_model (CrossValidatorModel): Объект CrossValidatorModel,
            содержащий лучшую модель.

    Returns:
        dict[str, float]: Параметры лучшей модели.
    """
    best_model = cv_model.bestModel
    best_params = {
        "regParam": best_model.stages[-1].getRegParam(),
        "elasticNetParam": best_model.stages[-1].getElasticNetParam(),
        "maxIter": best_model.stages[-1].getMaxIter()
    }
    return best_params

for key, value in get_best_model_params(cv_model=cv_model).items():
    print(f"{key}: {value}")

cv_scores = pd.DataFrame({
    "regParam": [grid[pipeline.getStages()[-1].regParam] for grid in param_grid],
    "elasticNetParam": [grid[pipeline.getStages()[-1].elasticNetParam]
                        for grid in param_grid],
    "cv_rmse": cv_model.avgMetrics,
})
print(cv_scores.sort_values("cv_rmse").to_string(index=False))

lr_model = cv_model.bestModel.stages[-1]
coefficients = np.array(lr_model.coefficients.toArray())

print(f"Размерность вектора признаков: {coefficients.size}")
print(f"Обнулено регуляризацией коэффициентов: {(coefficients == 0).sum()}")
print(f"Наибольший по модулю коэффициент: {np.abs(coefficients).max():.4f}")
print(f"Свободный член (intercept): {lr_model.intercept:.2f}")

def plot_training_summary(cv_model: CrossValidatorModel) -> None:
    """
    Отображает график зависимости значения ошибки от номера итерации на
    обучающей выборке, а также значения метрик RMSE и R^2.

    Args:
        cv_model (CrossValidatorModel): Обученная модель с использованием
            кросс-валидации.
    """
    # Получаем лучшую модель
    best_model = cv_model.bestModel

    # Получаем информацию о процессе обучения
    training_summary = best_model.stages[-1].summary

    # Получаем значения ошибки для каждой итерации
    objective_history = training_summary.objectiveHistory

    # Строим график зависимости значения ошибки от номера итерации
    plt.figure(figsize=(10, 6))
    sns.lineplot(x=range(len(objective_history)),
                 y=objective_history,
                 marker='o')
    plt.xlabel('Итерация')
    plt.ylabel('Ошибка')
    plt.title("Зависимость значения функции ошибки от номера итерации")

    # Получаем значения метрик
    rmse = training_summary.rootMeanSquaredError
    r2 = training_summary.r2

    # Добавляем значения метрик на график
    plt.text(0.95, 0.95, f"RMSE: {rmse:.2f}\nR^2: {r2:.4f}",
             transform=plt.gca().transAxes, ha='right', va='top',
             bbox=dict(facecolor='white', alpha=0.8), zorder=5)
    plt.grid()

    plt.show()

plot_training_summary(cv_model)

# Получаем датасет предсказаний
test_df_predictions = cv_model.transform(test_df)

# Извлекаем список колонок, устанавливаем метку и предсказание на последнее место
right_columns_order = test_df_predictions.columns
right_columns_order.remove(label_col)
right_columns_order.append(label_col)
right_columns_order.remove("prediction")
right_columns_order.append("prediction")

# Изменяем последовательность колонок и выводим датафрейм
test_df_predictions = test_df_predictions.select(*right_columns_order)
test_df_predictions.show()

def evaluate_model(data: DataFrame, metric_name: str) -> float:
    """
    Оценивает модель с использованием указанной метрики.

    Args:
        data (DataFrame): DataFrame, содержащий предсказания и фактические метки.
        metric_name (str): Название метрики для оценки
            (например, "rmse", "r2", "mae").

    Returns:
        float: Значение указанной метрики.
    """
    evaluator = RegressionEvaluator(labelCol=label_col,
                                    predictionCol="prediction",
                                    metricName=metric_name)
    metric_value = evaluator.evaluate(data)
    return metric_value

test_rmse = evaluate_model(test_df_predictions, "rmse")
test_r2 = evaluate_model(test_df_predictions, "r2")
test_mae = evaluate_model(test_df_predictions, "mae")

print(f"RMSE on test data: {test_rmse:.2f}")
print(f"R^2  on test data: {test_r2:.4f}")
print(f"MAE  on test data: {test_mae:.2f}")

median_duration = train_df.approxQuantile(label_col, [0.5], 0.001)[0]
print(f"Медиана Duration_min в обучающей выборке: {median_duration:.2f} мин")

baseline_rmse = (
    test_df
    .withColumn("residual", col(label_col) - median_duration)
    .select(sqrt(mean(col("residual") * col("residual"))).alias("rmse"))
    .collect()[0]["rmse"]
)
print(f"Baseline RMSE (константа = медиана): {baseline_rmse:.2f}")
print(f"Изменение RMSE относительно baseline: "
      f"{100.0 * (test_rmse - baseline_rmse) / baseline_rmse:.2f} %")


predictions_pandas = (
    test_df_predictions
    .sample(withReplacement=False, fraction=0.02, seed=17)
    .select(label_col, "prediction")
    .toPandas()
)
predictions_pandas["bucket"] = pd.qcut(predictions_pandas[label_col],
                                       q=10,
                                       labels=False,
                                       duplicates="drop")

bucketed = (
    predictions_pandas
    .groupby("bucket")
    .agg(actual=(label_col, "mean"),
         predicted=("prediction", "mean"),
         n=(label_col, "size"))
)
print(bucketed.round(2).to_string())

student_hdfs_folder = "otpushchennikov_directory"
# Получаем имя пользователя
user_name = os.getenv("USER")

# Путь модели в HDFS
model_hdfs_path = f"hdfs:///user/{user_name}/{student_hdfs_folder}/models/lab2-lr-duration-model"

# Сохраняем модель конвейера в HDFS
try:
    cv_model.bestModel.save(model_hdfs_path)
    print(f"Модель успешно сохранена в \"{model_hdfs_path}\"")
except Exception as e:
    print(f"Ошибка при сохранении модели: {e}")

spark.catalog.clearCache()
spark.stop()
\end{code}

\chapter{Листинг программного кода лабораторной работы №2 второй части}\label{app:D}\vspace{\baselineskip}

\begin{code}
import os
import time
import numpy as np
import pandas as pd
import seaborn as sns
import matplotlib.pyplot as plt
from pyspark.sql import SparkSession, DataFrame, Window
from pyspark import SparkConf
from pyspark.sql.functions import (
    col, lit, coalesce, when, desc, avg, sum as spark_sum,
    array_repeat, explode, approx_count_distinct
)
from pyspark.sql.types import DoubleType, IntegerType
from pyspark.ml import Pipeline
from pyspark.ml.feature import (
    VectorAssembler,
    StringIndexer,
    Binarizer,
    Bucketizer
)
from pyspark.ml.functions import vector_to_array
from pyspark.ml.classification import (
    RandomForestClassifier,
    RandomForestClassificationModel
)
from pyspark.ml.evaluation import BinaryClassificationEvaluator
from pyspark.ml.tuning import CrossValidator, CrossValidatorModel, ParamGridBuilder

def create_spark_configuration() -> SparkConf:
    """
    Создает и конфигурирует экземпляр SparkConf для приложения Spark.

    Returns:
        SparkConf: Настроенный экземпляр SparkConf.
    """
    # Получаем имя пользователя
    user_name = os.getenv("USER")

    conf = SparkConf()
    conf.setAppName("lab 2 US Accidents. Random Forest Classification")
    conf.setMaster("yarn")
    conf.set("spark.submit.deployMode", "client")
    conf.set("spark.executor.memory", "12g")
    conf.set("spark.executor.cores", "8")
    conf.set("spark.executor.instances", "2")
    conf.set("spark.driver.memory", "4g")
    conf.set("spark.driver.cores", "2")
    conf.set("spark.jars.packages",
             "org.apache.iceberg:iceberg-spark-runtime-3.5_2.12:1.6.0")
    conf.set("spark.sql.extensions",
             "org.apache.iceberg.spark.extensions.IcebergSparkSessionExtensions")
    conf.set("spark.sql.catalog.spark_catalog",
             "org.apache.iceberg.spark.SparkCatalog")
    conf.set("spark.sql.catalog.spark_catalog.type", "hadoop")
    conf.set("spark.sql.catalog.spark_catalog.warehouse",
             f"hdfs:///user/{user_name}/warehouse")
    conf.set("spark.sql.catalog.spark_catalog.io-impl",
             "org.apache.iceberg.hadoop.HadoopFileIO")

    return conf

conf = create_spark_configuration()

spark = SparkSession.builder.config(conf=conf).getOrCreate()
spark

database_name = "otpushchennikov_database"

spark.catalog.setCurrentDatabase(database_name)

df = spark.table("otpushchennikov_us_accidents_processed")

df.show()

df.printSchema()

df.count()

binary_flags = [
    "Crossing", "Give_Way", "Junction", "No_Exit",
    "Railway", "Station", "Stop", "Traffic_Signal"
]
float_features = [
    "Distance(mi)", "Temperature(F)", "Wind_Chill(F)", "Humidity(%)",
    "Pressure(in)", "Visibility(mi)", "Wind_Speed(mph)", "Precipitation(in)",
    "Month", "Day_Of_Week", "Hour"
]

label_col = "is_severe"

rows_before = df.count()
df = df.filter(col("Severity").isNotNull())
print(f"Строк без значения Severity: {rows_before - df.count()}")

df = df.withColumn(label_col,
                   when(col("Severity") >= 3, 1).otherwise(0).cast(IntegerType()))
df = df.withColumn("Duration_min", col("Duration_min").cast(DoubleType()))
for name in float_features:
    df = df.withColumn(name, col(name).cast(DoubleType()))
for flag in binary_flags:
    df = df.withColumn(
        flag,
        coalesce(col(flag).cast(DoubleType()), lit(0.0))
    )

df.groupBy(label_col).count().orderBy(desc(label_col)).show()

def stratified_train_test_split(data: DataFrame,
                                label_col: str,
                                ratio: float) -> tuple[DataFrame, DataFrame]:
    """
    Разделяет DataFrame на тренировочный и тестовый наборы с учетом стратификации.

    Args:
        data: Исходный DataFrame.
        label_col: Название колонки с меткой.
        ratio: Пропорция разделения данных.

    Returns:
        tuple[DataFrame, DataFrame]: Кортеж из тренировочного и тестового DataFrame.
    """
    # Проверяем корректность доли разделения
    assert (isinstance(ratio, float) and (0.0 <= ratio <= 1.0))

    # Формируем разделение для положительных и отрицательных объектов раздельно
    train_df_pos, test_df_pos = (data
                                 .filter(col(label_col) == 1)
                                 .randomSplit([ratio, 1 - ratio]))
    train_df_neg, test_df_neg = (data
                                 .filter(col(label_col) == 0)
                                 .randomSplit([ratio, 1 - ratio]))

    # Объединяем датафреймы
    return (train_df_pos.union(train_df_neg),
            test_df_pos.union(test_df_neg))

train_df, test_df = stratified_train_test_split(df, label_col, 0.8)

def oversample(data: DataFrame, column: str) -> DataFrame:
    """
    Выполняет oversampling положительных классов в DataFrame.

    Args:
        data: Исходный DataFrame.
        column: Название столбца с меткой.

    Returns:
        DataFrame: Датафрейм с выполненным oversampling.
    """
    # Разделим датафрейм на положительные и отрицательные классы
    pos = data.filter(col(column) == 1)
    neg = data.filter(col(column) == 0)

    # Получим количество записей в каждом классе
    total_pos = pos.count()
    total_neg = neg.count()

    # Если количество положительных классов меньше отрицательных,
    # выполним oversampling
    if total_pos < total_neg:
        # Вычислим количество необходимых дубликатов
        num_duplicates = total_neg - total_pos

        # Дублируем положительные записи: каждая строка разворачивается
        # в требуемое число копий
        oversampled_pos = (
            pos
            .withColumn("dummy", array_repeat(lit(1),
                                              num_duplicates // total_pos + 1))
            .withColumn("dummy", explode("dummy"))
            .drop("dummy")
        )

        # Объединим дублированные положительные записи с отрицательными
        balanced_df = neg.union(oversampled_pos)
    else:
        balanced_df = data

    return balanced_df

train_df = oversample(train_df, column=label_col)

train_df.groupBy(label_col).count().orderBy(desc(label_col)).show()

train_df = train_df.cache()
test_df = test_df.cache()

print(f"Train dataset size: {train_df.count()}")
print(f"Test  dataset size: {test_df.count()}")

categorical_features = ["State", "Weather_Condition",
                      "Wind_Direction", "Sunrise_Sunset"]
binary_features = [
    "Crossing", "Give_Way", "Junction", "No_Exit",
    "Railway", "Station", "Stop", "Traffic_Signal"
]
numeric_features = [
    "Distance(mi)", "Temperature(F)", "Wind_Chill(F)", "Humidity(%)",
    "Pressure(in)", "Visibility(mi)", "Wind_Speed(mph)", "Duration_min",
    "Month", "Day_Of_Week", "Hour"
]
binarizable_feature = "Precipitation(in)"
threshold = 0.1

cardinality = train_df.agg(
    *[approx_count_distinct(col(name)).alias(name) for name in categorical_features]
).toPandas().T
cardinality.columns = ["число уровней"]
print(cardinality.sort_values("число уровней", ascending=False).to_string())

def create_pipeline(categorical_features: list[str], numeric_features: list[str],
                    binary_features: list[str], binarized_col: str,
                    threshold: float, label_col: str, num_trees: int,
                    max_bins: int, seed: int) -> Pipeline:
    """
    Создает конвейер подготовки признаков и модель случайного леса.

    Args:
        categorical_features: Список категориальных признаков.
        numeric_features: Список количественных признаков.
        binary_features: Список бинарных признаков.
        binarized_col: Название бинаризуемого признака.
        threshold: Порог бинаризации.
        label_col: Название столбца с меткой.
        num_trees: Число деревьев в лесу.
        max_bins: Число корзин: должно быть не меньше числа уровней
            самого «широкого» категориального признака.
        seed: Зерно генератора случайных чисел (для воспроизводимости).

    Returns:
        Pipeline: Конвейер подготовки признаков и обучения модели.
    """
    # Формируем названия колонок для преобразованных признаков
    indexed_categorical_features = [f"{feature}_index" for feature in categorical_features]

    # Выполняем преобразования данных
    string_indexer = StringIndexer(inputCols=categorical_features,
                                   outputCols=indexed_categorical_features,
                                   handleInvalid="keep")
    binarizer = Binarizer(inputCol=binarized_col,
                          outputCol="precipitation_binary",
                          threshold=threshold)
    vector_all_assembler = VectorAssembler(
        inputCols=(indexed_categorical_features + binary_features +
                   numeric_features + ["precipitation_binary"]),
        outputCol="features",
        handleInvalid="keep"
    )

    # Создаем модель случайного леса
    random_forest = RandomForestClassifier(featuresCol="features",
                                           labelCol=label_col,
                                           predictionCol="prediction",
                                           numTrees=num_trees,
                                           maxBins=max_bins,
                                           seed=seed)

    # Создаем конвейер
    pipeline = Pipeline(stages=[
        string_indexer, binarizer, vector_all_assembler, random_forest
    ])

    return pipeline

pipeline = create_pipeline(categorical_features=categorical_features,
                           numeric_features=numeric_features,
                           binary_features=binary_features,
                           binarized_col=binarizable_feature,
                           threshold=threshold,
                           label_col=label_col,
                           num_trees=50,
                           max_bins=64,
                           seed=2026)

feature_names = (
    [f"{feature}_index" for feature in categorical_features]
    + binary_features
    + numeric_features
    + ["precipitation_binary"]
)

t0 = time.time()
single_pipeline_model = pipeline.fit(train_df)
one_fit_seconds = time.time() - t0

print(f"Время одного обучения конвейера со случайным лесом: {one_fit_seconds:.1f} сек")
print(f"Оценка времени полной сетки 3 x 3 при пяти фолдах: "
      f"{one_fit_seconds * 45 / 3600:.1f} ч")

# param_grid = ParamGridBuilder() \
#     .addGrid(pipeline.getStages()[-1].maxDepth, [4, 6, 8]) \
#     .addGrid(pipeline.getStages()[-1].numTrees, [50, 100, 200]) \
#     .build()

param_grid = ParamGridBuilder() \
    .addGrid(pipeline.getStages()[-1].maxDepth, [4, 6]) \
    .addGrid(pipeline.getStages()[-1].numTrees, [50]) \
    .build()
print(f"Число комбинаций параметров: {len(param_grid)}")

cv_evaluator = BinaryClassificationEvaluator(labelCol=label_col)

cross_validator = CrossValidator(estimator=pipeline,
                                 estimatorParamMaps=param_grid,
                                 evaluator=cv_evaluator,
                                 numFolds=5)

cv_model = cross_validator.fit(train_df)

def get_best_model_params(cv_model: CrossValidatorModel) -> dict[str, int]:
    """
    Получает параметры лучшей модели из объекта CrossValidatorModel.

    Args:
        cv_model: Объект CrossValidatorModel, содержащий лучшую модель.

    Returns:
        dict[str, int]: Параметры лучшей модели.
    """
    best_model = cv_model.bestModel
    best_params = {
        "maxDepth": best_model.stages[-1].getMaxDepth(),
        "numTrees": best_model.stages[-1].getNumTrees,
        "maxBins": best_model.stages[-1].getMaxBins()
    }
    return best_params

for key, value in get_best_model_params(cv_model=cv_model).items():
    print(f"{key}: {value}")

cv_scores = pd.DataFrame({
    "maxDepth": [grid[pipeline.getStages()[-1].maxDepth] for grid in param_grid],
    "numTrees": [grid[pipeline.getStages()[-1].numTrees] for grid in param_grid],
    "cv_auc_roc": cv_model.avgMetrics,
})
print(cv_scores.sort_values("cv_auc_roc", ascending=False).to_string(index=False))

rf_model = cv_model.bestModel.stages[-1]
importances = np.array(rf_model.featureImportances.toArray())

importance_df = (
    pd.DataFrame({"признак": feature_names, "важность": importances})
    .sort_values("важность", ascending=False)
)
print(importance_df.head(15).to_string(index=False))
print(f"Суммарная важность топ-5 признаков: {importance_df['важность'].head(5).sum():.4f}")

# Получаем датасет предсказаний
test_df_predictions = cv_model.transform(test_df)

# Извлекаем список колонок, устанавливаем метку и предсказание на последнее место
right_columns_order = test_df_predictions.columns
right_columns_order.remove(label_col)
right_columns_order.append(label_col)
right_columns_order.remove("prediction")
right_columns_order.append("prediction")

# Изменяем последовательность колонок и выводим датафрейм
test_df_predictions = test_df_predictions.select(*right_columns_order)
test_df_predictions.show()

def evaluate_model(data: DataFrame, label_col: str) -> dict[str, float]:
    """
    Оценивает модель с использованием метрик точности, полноты и F1-score.

    Args:
        data: DataFrame, содержащий предсказания и фактические метки.
        label_col: Название колонки с меткой.

    Returns:
        dict[str, float]: Словарь с метриками точности, полноты и F1-score.
    """
    # Вычисляем TP, FP, FN
    tp = data.filter((col(label_col) == 1) &
                     (col("prediction") == 1)).count()
    fp = data.filter((col(label_col) == 0) &
                     (col("prediction") == 1)).count()
    fn = data.filter((col(label_col) == 1) &
                     (col("prediction") == 0)).count()

    # Вычисляем метрики
    precision = tp / (tp + fp) if (tp + fp) > 0 else 0
    recall = tp / (tp + fn) if (tp + fn) > 0 else 0
    f1 = (2 * (precision * recall) / (precision + recall)
          if (precision + recall) > 0 else 0)

    # Возвращаем словарь с метриками
    return {
        "precision": precision,
        "recall": recall,
        "f1": f1
    }

metrics = evaluate_model(test_df_predictions, label_col)
print(f"Metrics: {metrics}")

counts = (
    test_df
    .agg(spark_sum(when(col(label_col) == 1, 1).otherwise(0)).alias("pos"),
         spark_sum(when(col(label_col) == 0, 1).otherwise(0)).alias("neg"))
    .first()
)
total_test = counts["pos"] + counts["neg"]

print(f"Доля мажоритарного класса (accuracy «всегда 0»): "
      f"{100.0 * counts['neg'] / total_test:.2f} %")
print(f"Доля миноритарного класса (precision «всегда 1»): "
      f"{100.0 * counts['pos'] / total_test:.2f} %")
    
def calculate_auc_roc(model: RandomForestClassificationModel,
                      test_data: DataFrame,
                      label_col: str) -> float:
    """
    Рассчитывает значение AUC ROC для модели.

    Args:
        model: Обученная модель классификации.
        test_data: Тестовый набор данных.
        label_col: Название столбца с меткой.

    Returns:
        float: Значение AUC ROC.
    """
    # Получаем прогнозы модели
    predictions = model.transform(test_data)

    # Создаем "оценщик" бинарной классификации
    evaluator = BinaryClassificationEvaluator(
        rawPredictionCol="rawPrediction",
        labelCol=label_col
    )

    # Вычисляем AUC ROC
    auc_roc = evaluator.evaluate(predictions)

    return auc_roc

def prepare_metrics_df(predictions: DataFrame, label_col: str) -> DataFrame:
    """
    Подготавливает DataFrame для расчета FPR, TPR и Precision.

    Args:
        predictions: DataFrame с прогнозами модели.
        label_col: Название колонки с меткой.

    Returns:
        DataFrame: DataFrame с рассчитанными столбцами FPR, TPR и Precision.
    """
    # Формируем датафрейм, содержащий результат предсказания,
    # вероятность предсказания и метку
    metrics_df = (
        predictions
        .select(
            vector_to_array(col("probability")).getItem(1).alias("probability"),
            col(label_col).alias("label"),
            col("prediction").cast("int").alias("prediction")
        )
    )

    # Сортируем датафрейм по значениям вероятности
    metrics_df = metrics_df.orderBy("probability")

    # Рассчитываем количество положительных и отрицательных классов
    pos_count = metrics_df.filter(metrics_df.label == 1).count()
    neg_count = metrics_df.filter(metrics_df.label == 0).count()

    # Рассчитываем значения FPR и TPR
    window = Window.partitionBy().orderBy(desc("probability"))
    metrics_df = (
        metrics_df
        .withColumn("FP_cum",
                    spark_sum(when(col("label") == 0, 1).otherwise(0)).over(window))
        .withColumn("TP_cum", spark_sum(col("label")).over(window))
        .withColumn("FPR", col("FP_cum") / neg_count)
        .withColumn("TPR", col("TP_cum") / pos_count)
        .withColumn("Precision",
                    col("TP_cum") / (col("TP_cum") + col("FP_cum")))
    )

    return metrics_df

def bucketize_metrics_df(metrics_df: DataFrame,
                         num_buckets: int = 1000) -> DataFrame:
    """
    Применяет Bucketizer к DataFrame для уменьшения объёма данных,
    возвращаемых на драйвер.

    Args:
        metrics_df: DataFrame с рассчитанными столбцами FPR, TPR и Precision.
        num_buckets: Количество бакетов.

    Returns:
        DataFrame: Бакетизированный DataFrame.
    """
    # Определяем границы бакетов. Верхнюю границу сдвигаем выше единицы,
    # иначе строки с вероятностью ровно 1,0 выходят за последнее значение
    # splits и Bucketizer с handleInvalid="error" аварийно остановит этап.
    splits = list(np.linspace(0, 1, num_buckets + 1))
    splits[-1] = 1.0 + 1e-9
    bucketizer = Bucketizer(
        splits=splits,
        inputCol="probability",
        outputCol="bucket"
    )

    # Применяем Bucketizer
    metrics_df = bucketizer.setHandleInvalid("error").transform(metrics_df)

    # Формируем одну точку из каждого бакета
    metrics_df = metrics_df.groupBy("bucket").agg(
        avg(col("probability")).alias("probability"),
        avg(col("FPR")).alias("FPR"),
        avg(col("TPR")).alias("TPR"),
        avg(col("Precision")).alias("Precision")
    )

    return metrics_df

def calculate_metrics(model: RandomForestClassificationModel,
                      test_data: DataFrame,
                      label_col: str) -> pd.DataFrame:
    """
    Рассчитывает значения FPR, TPR и Precision для построения ROC и PR-кривых.

    Args:
        model: Обученная модель классификации.
        test_data: Тестовый набор данных.
        label_col: Название столбца с меткой.

    Returns:
        pd.DataFrame: DataFrame со значениями FPR, TPR и Precision.
    """
    # Получаем прогнозы модели
    predictions = model.transform(test_data)

    # Подготавливаем DataFrame для расчета FPR, TPR и Precision
    metrics_df = prepare_metrics_df(predictions, label_col)

    # Применяем Bucketizer к DataFrame
    metrics_df = bucketize_metrics_df(metrics_df)

    return metrics_df.toPandas()

def plot_roc_pr_curves(pd_data: pd.DataFrame,
                       auc_roc: float) -> None:
    """
    Визуализирует ROC и PR-кривые, добавляет на график AUC ROC.

    Args:
        pd_data: Pandas DataFrame, содержащий столбцы "FPR", "TPR", "Precision".
        auc_roc: Значение AUC-ROC.

    Returns:
        None
    """
    fig, axes = plt.subplots(1, 2, figsize=(20, 6))

    # Построение ROC-кривой
    sns.lineplot(x="FPR", y="TPR", data=pd_data, ax=axes[0])
    sns.lineplot(x=[0, 1], y=[0, 1],
                 color="gray", linestyle="--", ax=axes[0])
    axes[0].set_xlabel("False Positive Rate (FPR)")
    axes[0].set_ylabel("True Positive Rate (TPR)")
    axes[0].set_title("ROC Curve")
    axes[0].grid()

    # Добавляем значение AUC-ROC в правом нижнем углу графика
    axes[0].text(0.94, 0.06, f'AUC-ROC: {auc_roc:.2f}',
                 transform=axes[0].transAxes,
                 fontsize=12, verticalalignment='bottom',
                 horizontalalignment='right',
                 bbox=dict(facecolor='white', alpha=0.8), zorder=5)

    # Построение PR-кривой
    sns.lineplot(x="TPR", y="Precision", data=pd_data, ax=axes[1])
    axes[1].set_xlabel("Recall")
    axes[1].set_ylabel("Precision")
    axes[1].set_title("PR Curve")
    axes[1].grid()

    plt.tight_layout()
    plt.show()

# Рассчитываем AUC ROC
auc_roc = calculate_auc_roc(model=cv_model,
                            test_data=test_df,
                            label_col=label_col)
pd_dataframe = calculate_metrics(model=cv_model,
                                 test_data=test_df,
                                 label_col=label_col)
print(f"AUC ROC: {auc_roc:.4f}")
plot_roc_pr_curves(pd_dataframe, auc_roc)

threshold_probability = pd_dataframe[pd_dataframe['TPR'] >= 0.60]['probability'].max()
print(f"Вероятность -- граница разделения, при которой TPR не меньше 60%: "
      f"{threshold_probability:.2f}")

cv_model.bestModel.stages[-1].setThresholds([1 - threshold_probability,
                                             threshold_probability])
test_df_predictions = cv_model.transform(test_df)
metrics = evaluate_model(test_df_predictions, label_col)
print(f"Metrics: {metrics}")

spark.catalog.clearCache()
spark.stop()
\end{code}